\begin{figure}[htbp]
\centering
\begin{tikzpicture}[
  x=1cm, y=1cm,
  % Двата панела делят една времева ос и една и съща сива лента за престоя.
  % Това е нарочно: сравнението има смисъл само ако престоят е един и същ.
  sbar/.style={draw, fill=gray!10, anchor=west, inner sep=0pt,
               minimum height=0.34cm},
  stalled/.style={sbar, fill=gray!40},
  % Един и същ знак за изпращане на заявка в двата панела. Броят му вътре в
  % лентата е цялото твърдение на фигурата, затова знакът трябва да е еднакъв.
  send/.style={line width=1pt},
  hdr/.style={anchor=west, font=\fontsize{9}{11}\selectfont},
  lane/.style={note, anchor=east, align=right},
  meas/.style={Stealth-Stealth, semithick},
  inbar/.style={font=\fontsize{7.5}{9}\selectfont},
]

% ----------------------------------------------------- престой на сървъра
% Лентата е прекъсната само там, където минава заглавието на долния панел.
% Двете парчета са на едни и същи x, за да се чете като един и същ престой.
\begin{scope}[on background layer]
  \fill[gray!12] (5.80,1.25) rectangle (8.00,-2.10);
  \fill[gray!12] (5.80,-2.78) rectangle (8.00,-6.00);
\end{scope}
\foreach \x in {5.80,8.00} {
  \draw[densely dashed, black!55] (\x,1.25) -- (\x,-2.10);
  \draw[densely dashed, black!55] (\x,-2.78) -- (\x,-6.00);
}
\node[note, anchor=south, align=center] at (6.90,1.32) {престой на сървъра};

% ===================================================== ЗАТВОРЕН ЦИКЪЛ
\node[hdr] at (0.35,1.95)
  {\textbf{Затворен цикъл} (\code{h2load}): връзката чака отговор преди следващата заявка.};

% Ред с изпращанията от трите връзки, слети на една ос.
\node[lane] at (1.85,1.02) {изпращания};
\foreach \x in {2.00,2.20,2.45,2.90,3.15,3.25,3.80,3.90,4.30,4.60,4.70,5.30,5.45,5.60}
  \draw[send] (\x,0.87) -- (\x,1.17);
\foreach \x in {8.15,8.20,8.25,8.90,9.15,9.25,9.65,10.05,10.30,10.40,10.95,11.15,11.35}
  \draw[send] (\x,0.87) -- (\x,1.17);

% Празнината между последното и първото следващо изпращане е цялата поанта.
\draw[meas] (5.60,0.70) -- (8.15,0.70);
\node[note, anchor=north, align=center] at (6.88,0.62) {нула изпращания};

% Три връзки. Всяка има точно една заявка в полет, когато престоят започне,
% затова всяка има точно една удължена лента и нищо друго вътре в престоя.
\node[lane] at (1.85,0.02) {връзка 1};
\foreach \x in {2.00,2.90,3.80,4.70} \node[sbar, minimum width=0.80cm] at (\x,0.02) {};
\node[stalled, minimum width=2.55cm] at (5.60,0.02) {};
\foreach \x in {8.25,9.15,10.05,10.95} \node[sbar, minimum width=0.80cm] at (\x,0.02) {};

\node[lane] at (1.85,-0.58) {връзка 2};
\foreach \x in {2.20,3.25,4.30} \node[sbar, minimum width=0.90cm] at (\x,-0.58) {};
\node[stalled, minimum width=2.65cm] at (5.45,-0.58) {};
\foreach \x in {8.20,9.25,10.30,11.35} \node[sbar, minimum width=0.90cm] at (\x,-0.58) {};

\node[lane] at (1.85,-1.18) {връзка 3};
\foreach \x in {2.45,3.15,3.90,4.60} \node[sbar, minimum width=0.60cm] at (\x,-1.18) {};
\node[stalled, minimum width=2.75cm] at (5.30,-1.18) {};
\foreach \x in {8.15,8.90,9.65,10.40,11.15} \node[sbar, minimum width=0.60cm] at (\x,-1.18) {};

% Мярката, която затвореният цикъл наистина дава.
\draw[densely dotted, gray!70] (5.30,-1.36) -- (5.30,-1.62);
\draw[densely dotted, gray!70] (8.05,-1.36) -- (8.05,-1.62);
\draw[meas] (5.30,-1.68) -- (8.05,-1.68);
\node[note, anchor=north, align=center, fill=white, inner sep=1.5pt] at (6.68,-1.76) {време за обслужване};

% ===================================================== ОТВОРЕН ЦИКЪЛ
\node[hdr] at (0.35,-2.50)
  {\textbf{Отворен цикъл} (\code{wrk2}): фиксиран темп на заявките, независимо от сървъра.};

\node[lane] at (1.85,-3.10) {заявки};
\draw[gray!55, thin] (1.95,-3.10) -- (12.30,-3.10);
\foreach \x in {2.00,2.43,...,12.30} \draw[send] (\x,-3.25) -- (\x,-2.95);

% Една проследена заявка. Пристига вътре в престоя, чака, после се обслужва.
\node[lane] at (1.85,-3.75) {една заявка};
\node[sbar, minimum width=2.60cm, inbar] at (6.30,-3.75) {чакане в опашката};
\node[stalled, minimum width=0.80cm] at (8.90,-3.75) {};
\node[note, anchor=west] at (9.80,-3.75) {обслужване};
\draw[densely dotted, gray!70] (6.30,-3.25) -- (6.30,-3.58);

\draw[meas] (6.30,-4.17) -- (9.70,-4.17);
\node[note, anchor=north, align=center, fill=white, inner sep=1.5pt] at (8.00,-4.25) {време за отговор};

% Дължина на опашката. Стълбата е стъпаловидна, защото пристиганията са
% дискретни: по едно качване на всяко пристигане, докато сървърът стои.
\node[lane] at (1.85,-5.30) {опашка};
\fill[gray!25]
  (2.00,-5.90) -- (5.87,-5.90) -- (5.87,-5.70) -- (6.30,-5.70) -- (6.30,-5.50) --
  (6.73,-5.50) -- (6.73,-5.30) -- (7.16,-5.30) -- (7.16,-5.10) -- (7.59,-5.10) --
  (7.59,-4.90) -- (8.02,-4.90) -- (8.02,-5.10) -- (8.45,-5.10) -- (8.45,-5.30) --
  (8.88,-5.30) -- (8.88,-5.50) -- (9.31,-5.50) -- (9.31,-5.70) -- (9.74,-5.70) --
  (9.74,-5.90) -- (12.30,-5.90) -- cycle;
\draw[thick]
  (2.00,-5.90) -- (5.87,-5.90) -- (5.87,-5.70) -- (6.30,-5.70) -- (6.30,-5.50) --
  (6.73,-5.50) -- (6.73,-5.30) -- (7.16,-5.30) -- (7.16,-5.10) -- (7.59,-5.10) --
  (7.59,-4.90) -- (8.02,-4.90) -- (8.02,-5.10) -- (8.45,-5.10) -- (8.45,-5.30) --
  (8.88,-5.30) -- (8.88,-5.50) -- (9.31,-5.50) -- (9.31,-5.70) -- (9.74,-5.70) --
  (9.74,-5.90) -- (12.30,-5.90);

% Основата на опашката служи и за времева ос на цялата фигура.
\draw[-Stealth, semithick] (1.95,-5.90) -- (12.55,-5.90);
\node[note, anchor=south east] at (12.55,-5.87) {време};

% ===================================================== следствие
\node[note, anchor=north west, text width=11.90cm] at (0.35,-6.30)
  {\textbf{Следствие.} p99 от затворен цикъл подценява опашката на разпределението.
   Клиентът е спрял да поражда точно онова натоварване, което би я произвело.};

\end{tikzpicture}
\caption{Съгласувано пропускане, показано като механизъм, а не като твърдение.
Горе всяка от $C$ връзки изпраща следваща заявка едва след получен отговор, затова
за целия престой на сървъра няма нито едно изпращане. Измерва се време за обслужване
на допуснатите заявки. Долу заявките се предлагат с постоянен темп, престоят не ги
спира, опашката расте и се източва след него. Измерва се време за отговор, тоест
чакане плюс обслужване. Затова 99-ият персентил от затворен цикъл подценява горния
край на разпределението: липсват именно заявките, които престоят би забавил най-много.}
\label{fig:closed_vs_open_loop}
\end{figure}
